\documentclass[12pt,a4paper]{article}
\usepackage{amsmath,amssymb,amsthm}
\usepackage{hyperref}
\usepackage{geometry}
\usepackage{enumitem}
\geometry{margin=2.5cm}

\newtheorem{theorem}{Theorem}[section]
\newtheorem{lemma}[theorem]{Lemma}
\newtheorem{corollary}[theorem]{Corollary}
\theoremstyle{definition}
\newtheorem{definition}[theorem]{Definition}
\theoremstyle{remark}
\newtheorem{remark}[theorem]{Remark}

\title{Protein Folding from Recognition Science:\\
Native Fold as $J$-Cost Minimum on the $\varphi$-Ladder}

\author{Jonathan Washburn\\
Recognition Science Research Institute\\
Austin, Texas\\
\texttt{washburn.jonathan@gmail.com}}

\date{March 2026}

\begin{document}
\maketitle

\begin{abstract}
We derive the protein folding problem from Recognition Science (RS).
The native fold is the global minimum of the recognition energy $Q_6$,
defined as the sum of $J$-costs over 8-tick windows across the
conformational landscape, where $J(x) = \frac{1}{2}(x + x^{-1}) - 1$.
The folding funnel maps to the $J$-cost landscape: the native
structure is the unique, stable global minimum.  Levinthal's paradox
is resolved: the $\varphi$-ladder provides a hierarchical folding
pathway that guides the search, so the system need not sample all
configurations.  The native fold is stable (energy gap $\geq c_{\min}
\cdot \mathrm{defect}$) and unique (equal $Q_6$ implies equivalent
conformations).
\end{abstract}

%% ============================================================
\section{Introduction}
\label{sec:intro}
%% ============================================================

Anfinsen's dogma~\cite{Anfinsen} states that a protein's native fold is
determined by its amino acid sequence: the native structure is the
thermodynamic ground state.  The physical basis of this hypothesis
remains incompletely understood.  A central puzzle is Levinthal's
paradox~\cite{Levinthal}: a typical protein has $\sim 10^{100}$ or more
possible conformations; if folding required random search, the time to
find the native state would exceed the age of the universe.  Yet
proteins fold in milliseconds.  The resolution requires that the
energy landscape be funnel-like, guiding the chain toward the native
state rather than requiring exhaustive search.

Recognition Science provides a principled derivation.  The cost
functional $J(x) = \frac{1}{2}(x + x^{-1}) - 1$ and the $\varphi$-ladder
($\varphi = (1+\sqrt{5})/2$, $E_{\mathrm{coh}} = \varphi^{-5}~\mathrm{eV}$)
define a recognition energy landscape.  We show that the native fold
minimizes this landscape and that the $\varphi$-ladder structure
resolves Levinthal's paradox.

%% ============================================================
\section{The $J$-Cost and Recognition Energy}
\label{sec:jcost}
%% ============================================================

\begin{definition}[$J$-Cost]
\label{def:jcost}
The recognition cost functional is
\begin{equation}
  J(x) = \frac{1}{2}(x + x^{-1}) - 1, \qquad x > 0.
\end{equation}
$J(x) \geq 0$ with equality iff $x = 1$; $J$ penalizes deviation from
$\varphi$-resonance.
\end{definition}

\begin{definition}[$Q_6$ Recognition Energy]
\label{def:q6}
The recognition energy of a conformation $c$ is
\begin{equation}
  Q_6(c) = \sum_{\text{windows}~w} \sum_k J(w_k),
\end{equation}
where $w$ ranges over 8-tick windows ($2^D = 8$ for $D = 3$) and
$J(w_k)$ is the cost of each window component $w_k$.
\end{definition}

\begin{definition}[Recognition Defect]
\label{def:defect}
$\mathrm{defect}(\mathrm{native}, c) = Q_6(c) - Q_6(\mathrm{native})$.
\end{definition}

%% ============================================================
\section{Folding Funnel as $J$-Cost Landscape}
\label{sec:funnel}
%% ============================================================

The empirical folding funnel---a landscape that narrows toward the
native state---maps directly to the $J$-cost landscape.  High-energy
(misfolded) conformations have large $Q_6$; as the chain approaches
$\varphi$-resonant geometry, $Q_6$ decreases.  The native fold is the
global minimum.

\begin{theorem}[Global Minimum]
\label{thm:global}
The native fold minimizes $Q_6$:
\begin{equation}
  Q_6(\mathrm{native}) \leq Q_6(c) \quad \text{for all conformations } c.
\end{equation}
\end{theorem}

\begin{theorem}[Stability]
\label{thm:stability}
The native fold is stable with energy gap
\begin{equation}
  Q_6(c) - Q_6(\mathrm{native}) \geq c_{\min} \cdot
  \mathrm{defect}(\mathrm{native}, c).
\end{equation}
\end{theorem}

\begin{theorem}[Uniqueness]
\label{thm:uniqueness}
If $Q_6(c_1) = Q_6(c_2) = Q_6(\mathrm{native})$, then
$\mathrm{defect} = 0$, so $c_1$ and $c_2$ are equivalent.
\end{theorem}

%% ============================================================
\section{Levinthal's Paradox Resolved}
\label{sec:levinthal}
%% ============================================================

\begin{theorem}[Hierarchical Folding Pathway]
\label{thm:levinthal}
The $\varphi$-ladder provides a hierarchical folding pathway.  The
system need not sample all configurations; the discrete rung structure
guides the search along a restricted set of $\varphi$-resonant
intermediate states.
\end{theorem}

\begin{proof}[Proof sketch]
The 8-tick epoch ($D = 3$) defines a coarse-grained partition of
conformational space.  Each 8-tick window has a finite number of
$J$-cost values determined by the $\varphi$-ladder.  The folding
dynamics proceed by local $J$-cost minimization within each window,
then coordination between windows.  The hierarchy (8-tick $\to$ longer
scales) constrains the search to a fraction of the full configuration
space.  The funnel shape emerges from the monotonic decrease of $J$
toward resonance.
\end{proof}

\begin{corollary}[Levinthal Resolution]
Levinthal's paradox assumes random search over all conformations.  RS
predicts that the $\varphi$-ladder restricts the effective search space
to a hierarchical pathway.  Folding time scales with the number of
rungs traversed, not with the total configuration count.
\end{corollary}

\begin{remark}
The 8-tick window size ($2^3 = 8$) is fixed by $D = 3$.  This
discrete structure is the origin of the funnel: the landscape is not
featureless but has built-in guidance toward the native state.
\end{remark}

%% ============================================================
\section{Predictions and Falsifiers}
\label{sec:falsifiers}
%% ============================================================

\begin{enumerate}[label=(\roman*)]
\item \textbf{Native fold as global minimum}: The native structure
  minimizes $Q_6$.  \textit{Falsifier}: Discovery of a conformation
  with lower recognition energy than the native fold for a
  well-characterized protein.

\item \textbf{Funnel landscape}: The $J$-cost landscape should be
  funnel-like (narrowing toward native).  \textit{Falsifier}: Folding
  kinetics inconsistent with funnel dynamics (e.g., multiple deep
  minima with comparable populations).

\item \textbf{Levinthal resolution}: Folding pathways should exhibit
  hierarchical structure consistent with 8-tick windows.
  \textit{Falsifier}: Folding that requires exhaustive search over
  $\varphi$-resonant states.

\item \textbf{Stability gap}: The energy gap between native and
  misfolded states should scale with defect.  \textit{Falsifier}:
  Proteins with negligible stability gap that nonetheless fold
  uniquely.
\end{enumerate}

%% ============================================================
\section{Conclusion}
\label{sec:conclusion}
%% ============================================================

Protein folding in RS is $J$-cost minimization on the $\varphi$-ladder.
The native fold is the unique, stable global minimum of $Q_6$.  The
folding funnel maps to the $J$-cost landscape, and Levinthal's paradox
is resolved by the hierarchical structure of the 8-tick epoch: the
$\varphi$-ladder guides the search along a restricted pathway rather
than requiring random exploration of all conformations.

\begin{thebibliography}{9}
\bibitem{Anfinsen}
C.~B.~Anfinsen, Science \textbf{181}, 223 (1973).

\bibitem{Levinthal}
C.~Levinthal, in ``Mossbauer Spectroscopy in Biological Systems''
(University of Illinois Press, 1969), p.~22.
\end{thebibliography}

\end{document}
